\section{Discussion and Operational Implications}
The research questions lead to a qualified but coherent conclusion. For RQ1, the common retrospective attributes support useful cancellation ranking within the random benchmark, and nonlinear models improve on the linear baseline in that population. However, repeated profiles and the independently tuned time check reveal a narrower scope than the random score suggests. For RQ2, forest scores are strongly associated with booking context, lead time, and historical cancellations. For RQ3, hotel type and the five numerical segments show distinct observed cancellation rates. For RQ4, a ranking score needs calibration assessment and an explicit decision policy before it can become an operational trigger.

The most immediate use case is prioritising a small set of bookings for confirmation or staff review. Longer-lead segments can motivate earlier reminders, while long-stay segments can inform capacity review. These are candidate actions to test, not measured benefits from this study. We have neither intervention outcomes nor reliable booking-level cancellation costs, so we cannot quantify recovered revenue, recommend an optimal deposit amount, or claim that a policy reduces cancellations.

With calibrated probabilities and known error costs, a simplified review rule would flag a booking when \(p>C_{FP}/(C_{FP}+C_{FN})\), assuming correct decisions have zero cost and the stated false-positive and false-negative costs apply. Actual review, contact, room replacement, and guest experience costs complicate that rule. An F1-selected threshold is a reproducible statistical benchmark rather than a revenue-maximising policy. Calibration on historical random-training predictions also does not guarantee calibration after temporal drift \citep{sklearncalibration}.

Several limitations remain. The observations come from two hotels in 2015--2017, so transfer to another hotel or period requires new evaluation. The source is retrospective and lacks a complete initial-booking snapshot, which prevents a verified real-time prediction claim. Full-data EDA is descriptive, and shared folds are not group-disjoint. Hyperparameter and threshold summaries are post-selection, not nested validation. The independently tuned time experiment covers the forest only and does not establish a clean chronological ranking of all four algorithms. Changes in prevalence also affect AP and calibration.

Finally, the data do not identify causal drivers of cancellation or the effect of changing guest attributes. Country and agent may encode sensitive contextual patterns; automated adverse treatment would be unsupported. A future operational study should establish feature timestamps, use group-aware time-respecting selection, validate on newer records, measure review capacity and error costs, and evaluate reminders prospectively. Transparent checks on evaluation boundaries are necessary because leakage can otherwise produce misleading performance claims \citep{kapoor2023}.

\section{Conclusion}
We provide a common analytical record linking data cleaning, descriptive booking patterns, training-fitted segmentation, classifier comparison, forest explanation, and probability assessment. Cancellation rates differ substantially across hotels and booking segments. The random forest outperforms the logistic baseline in the retained random benchmark, while unseen-profile and separately tuned time results caution against extending that finding to unfamiliar bookings. Model ranking, a thresholded decision, and an explanation answer different questions. The defensible practical conclusion is to treat these models as retrospective candidates for locally validated decision support, with field availability, prospective performance, and intervention value still to be established.
